\section{Common input for zero-cycles with coefficients and splitting transfers}\label{app:review-cycle-transfer}
The object of this appendix is to prove Corollary~\ref{cor:review-small-cd-splitting-transfer} and Lemma~\ref{lem:review-p-rank-two-frobenius-norm}, which together give Proposition~\ref{prop:review-splitting-transfer-norm-iso}: over a two-dimensional local field, the reduced norm maps the subgroup generated by splitting transfers isomorphically onto $K_2(F)$.
Away from the characteristic, we compute zero-cycles with Milnor $K$-coefficients on products of Severi--Brauer varieties; at the characteristic prime, we use the Frobenius.

Throughout this appendix, for a smooth variety $X$ over a field $k$, $M(X)$ denotes its effective motive with integral coefficients, and $\operatorname{Hom}$ denotes morphisms in the category of effective motives.
For a central simple algebra $B/k$, the sheaf with transfers $\Z_B$ takes the value $\ind(B_{k(U)})\Z\subset\Z$ on a smooth integral variety $U/k$, and $\Z_B(q)=\Z_B\otimes\Z(q)$ denotes its Tate twist.
We make motivic calculations with $\Z_{(\ell)}$ coefficients for a fixed prime $\ell$ different from the characteristic; if $\operatorname{char}k=p>0$, then $p$ is invertible in these coefficients.
The Frobenius calculation in Lemma~\ref{lem:review-p-rank-two-frobenius-norm} uses integral coefficients.
Recall (Definition~\ref{def:paper-transfer}) that $I_F(D)=\sum_E\theta_EK_2(E)$, the sum over all finite splitting fields $E/F$, separable or not, and that $\Nrd_2\theta_E=N_{E/F}$.

\begin{definition}\label{def:review-rost-zero-cycles}
For a finite-type $k$-scheme $Z$ and $j\ge0$, let $Z_{(a)}$ denote the set of points whose closures have dimension $a$.
The low-degree terms of the Rost complex are
\[
 C_a(Z,K_j^M)=\bigoplus_{z\in Z_{(a)}}K_{j+a}^M(\kappa(z))\quad(a=0,1),
 \qquad A_0(Z,K_j^M)=\coker(C_1\xrightarrow{\partial}C_0).
\]
The differential is the sum of the residue maps on the normalizations of the closures of one-dimensional points, followed by norms on residue fields.
Proper pushforward is defined by the residue-field norm when the image has the same dimension and by zero when the dimension drops; it commutes with the differential.
In particular, when $Z$ is proper over $k$, the structure map gives the degree map $\deg_Z:A_0(Z,K_j^M)\to K_j^M(k)$ \cite{Rost}.
\end{definition}

\begin{lemma}
\label{lem:review-proper-cycle-pushforward}
Let $f:X\to Y$ be a proper morphism of finite-type $k$-schemes, and let $j\ge0$.
Suppose that, for every closed point $y\in Y_{(0)}$, the map
\[
 A_0(X_y,K_j^M)\xrightarrow{\deg}K_j^M(\kappa(y))
\]
is an isomorphism, and that $A_0(X_\eta,K_{j+1}^M)\to K_{j+1}^M(\kappa(\eta))$ is surjective for every $\eta\in Y_{(1)}$ (the fibers $X_y$ and $X_\eta$ being regarded as schemes over $\kappa(y)$, resp.\ $\kappa(\eta)$).
Then $f_*:A_0(X,K_j^M)\to A_0(Y,K_j^M)$ is an isomorphism.
The same holds when both the hypotheses on the fibers and the conclusion are taken with $\Z_{(\ell)}$ coefficients.
\end{lemma}
\begin{proof}
Surjectivity follows by lifting the coefficient at each closed point of a zero-cycle from the corresponding fiber.
For injectivity, take $z\in C_0(X,K_j^M)$ and write $f_*z=\partial b$.
By assumption, the coefficient of $b$ at each one-dimensional point $\eta$ lifts from $K_{j+1}^M$ coefficients supported on finitely many closed points of $X_\eta$.
We take their closures in $X$, with the same coefficients over their function fields, to obtain a chain $\widetilde b\in C_1(X,K_j^M)$.
These closures are curves that are proper and generically finite over $\overline{\{\eta\}}$.
We have $f_*\partial\widetilde b=\partial b$ by the compatibility of residues on the normalizations of these closures with field norms, that is, by the chain-map property of proper pushforward.

Thus $f_*(z-\partial\widetilde b)=0$ in $C_0(Y,K_j^M)$.
Each part of $z-\partial\widetilde b$ on a closed fiber has degree zero, so injectivity of the degree map on the closed fibers makes it the boundary of a chain in $C_1(X_y,K_j^M)$.
Pushing these chains forward along the closed immersions $X_y\to X$ and adding them shows that $z-\partial\widetilde b$, and hence $z$, is a boundary.
\end{proof}

\begin{lemma}
\label{lem:review-prime-degree-cycle}
Suppose that $\ell\ne\operatorname{char}k$, $j\ge1$, and $\cd_\ell k\le j+1$, and let $B/k$ be a central simple algebra of degree $\ell$.
Then the degree map
\[
 A_0(\SB(B),K_j^M)_{(\ell)}\xrightarrow{\sim}K_j^M(k)_{(\ell)}
\]
is an isomorphism.
\end{lemma}
\begin{proof}
All isomorphisms in this proof are understood after $\ell$-localization.
First assume that $k$ is perfect, and put $X=\SB(B)$.
Consider the coniveau spectral sequence computing $\CH^{\ell-1+j}(X,j)$; with $n=\dim X+j$, its $E_1$-terms are $E_1^{a,b}=\bigoplus_{x\in X^{(a)}}H^{b-a}(\kappa(x),\Z(n-a))$.
On the target diagonal, degree minus weight equals $\dim X-a$, and at the source of $d_r$ it equals $r-1$, so the terms in codimension $a<\ell-1$ and the differentials $d_r$ with $r\ge2$ entering the target term vanish because the cohomological degree over a field exceeds the weight.
Only $K_j^M$ at closed points and residues from $K_{j+1}^M$ at one-dimensional points remain, so duality \cite[Corollary C.7.3]{KahnLevine2010} and the cancellation theorem give
\[
 A_0(X,K_j^M)=\CH^{\ell-1+j}(X,j)
       =\operatorname{Hom}(\Z,M(X)(j)[j]).
\]

By the finite tower \cite[(8.1)]{KahnLevine2010}, for which $f_\ell M(X)=0$, the $s$th slice of $M(X)$ is $\Z_{B^{\otimes(s+1)}}(s)[2s]$ for $0\le s\le\ell-1$.
For $s>0$, the contributing group
\[
 H^{j+2s}(k,\Z_{B^{\otimes(s+1)}}(j+s))
\]
and the group one degree higher vanish because degree exceeds weight.
Hence $A_0(X,K_j^M)\simeq H^j(k,\Z_B(j))$; on the bottom layer the structure map is $\Z_B\hookrightarrow\Z$, so the degree map corresponds to the reduced norm induced by this inclusion \cite[Lemmas 5.4.2, 8.1.1(3), Theorems 6.5.3, 7.4.2]{KahnLevine2010}.
The norm residue isomorphism theorem supplies the Beilinson--Lichtenbaum hypothesis of the Kahn--Levine exact sequence in every required weight \cite{Voevodsky}.
The sequence makes the degree map injective with cokernel contained in $H^{j+2}_{\et}(k,\Z/\ell(j+1))$, which vanishes by the cohomological dimension hypothesis, so the degree map is an isomorphism \cite[Theorem 8.2.2(2)]{KahnLevine2010}.

If $k$ is imperfect of positive characteristic, pass to its perfect closure; the finite \'etale sites are equivalent, so $\cd_\ell$ is preserved.
Zero-cycles with coefficients and their relations are defined by finitely many data and therefore descend to a finite purely inseparable stage.
Both injectivity and surjectivity descend to $k$ because restriction followed by transfer is multiplication by the degree of this stage, a power of $p$ that is invertible after $\ell$-localization.
\end{proof}

\begin{proposition}
\label{prop:review-small-cd-cycle-degree}
Suppose that $\ell\ne\operatorname{char}k$, $j\ge1$, and $\cd_\ell k\le j+1$.
For every finite product $Z/k$ of Severi--Brauer varieties of central simple algebras, the degree map
\[
 \deg_Z:A_0(Z,K_j^M)_{(\ell)}\xrightarrow{\sim}K_j^M(k)_{(\ell)}
\]
is an isomorphism.
\end{proposition}
\begin{proof}
We prove the assertion simultaneously for all pairs $(k,j)$ satisfying the hypotheses.
For an $a$-dimensional point of a finite-type $k$-scheme $W$, the residue field $\kappa$ satisfies $\cd_\ell\kappa\le\cd_\ell k+a\le(j+a)+1$ \cite[II, \S4.2, Proposition 11]{SerreGalois}.
Thus, when we apply Lemma~\ref{lem:review-proper-cycle-pushforward} to a projection, the base fields of both types of fibers satisfy the bound, with coefficients $K_j^M$ at closed points and $K_{j+1}^M$ at one-dimensional points.
For a product of factors $\SB(B_i)$ with $\deg B_i=\ell$, apply Lemma~\ref{lem:review-prime-degree-cycle} fiberwise and induct on the number of factors to see that the degree map is an isomorphism.

Now consider $X=\SB(B)$ for an arbitrary central simple algebra $B/k$.
We may pass to a finite separable extension of degree prime to $\ell$ that splits the other primary Brauer components, and adjoin $\mu_\ell$, which requires an extension of degree dividing $\ell-1$.
Restriction and transfer descend the degree isomorphism to the original field.
Assume henceforth that $[B]$ has period $\ell^r$, and induct on $r$.
For $r=0$, use the projective bundle formula.
We express $\ell^{r-1}[B]$ as a sum of classes of symbol algebras of degree $\ell$, and let $Y$ be the product of their Severi--Brauer varieties \cite{MerkurjevSuslin1982,Voevodsky}.

Every fiber of the projection $f:X\times Y\to X$ is a product of factors $\SB(B_i)$ with $\deg B_i=\ell$, so the preceding calculation and Lemma~\ref{lem:review-proper-cycle-pushforward} show that $f_*$ is an isomorphism.
Over every residue field $\kappa(y)$ of $Y$, the point $y$ splits each symbol algebra factor, so $\ell^{r-1}[B]=0$ over $\kappa(y)$.
The induction on the period therefore applies to the fibers of $g:X\times Y\to Y$ over closed and one-dimensional points, with coefficients $j$ and $j+1$, respectively, so $g_*$ is an isomorphism by the same lemma.
Transitivity of proper pushforward gives $\deg_X=\deg_Y\circ g_*\circ f_*^{-1}$, which is an isomorphism because $\deg_Y$ already is.
Finally, we adjoin one arbitrary algebra factor at a time and use the same fiberwise argument to prove the assertion for general finite products.
\end{proof}

\begin{corollary}
\label{cor:review-small-cd-splitting-transfer}
Suppose that $\ell\ne\operatorname{char}F$ and $\cd_\ell F\le3$.
For every central division algebra $D$ over $F$,
\[
 \Nrd_2:I_F(D)_{(\ell)}\xrightarrow{\sim}K_2(F)_{(\ell)}
\]
is an isomorphism.
\end{corollary}
\begin{proof}
Put $X=\SB(D)$.
Compose the Morita equivalence given by the universal right ideal at each closed point with proper pushforward to obtain $\bigoplus_{t\in X_{(0)}}K_2^M(\kappa(t))\to K_2(D)$.
On the normalization of the closure of each one-dimensional point, the pullback of the universal ideal splits the Azumaya algebra.
There the composite of the map from Milnor $K_3$ to Quillen $K_3$ with the localization boundary and proper pushforward is zero, since consecutive maps in the localization sequence compose to zero.
The Rost boundary relations therefore vanish, giving a map $\Theta:A_0(X,K_2^M)\to K_2(D)$.
Every closed point $t$ splits $D$, and the component of $\Theta$ at $t$ is the splitting transfer $\theta_{\kappa(t)}$.
Conversely, a rational point of $X$ over a finite splitting field $E/F$ maps to a closed point $t$, and we have $\theta_E=\theta_{\kappa(t)}N_{E/\kappa(t)}$.
Thus $\im\Theta=I_F(D)$.
Since the reduced norm is compatible with the field norm at each closed point, we have $\Nrd_2\Theta=\deg_X$ \cite{Rost,QuillenHigher,KahnLevine2010}.
By Proposition~\ref{prop:review-small-cd-cycle-degree} with $j=2$, the degree map is an isomorphism after $\ell$-localization, so the reduced norm on the image $I_F(D)_{(\ell)}$ of $\Theta_{(\ell)}$ is an isomorphism.
\end{proof}

\begin{lemma}
\label{lem:review-p-rank-two-frobenius-norm}
Suppose that $\operatorname{char}F=p$ and $[F:F^p]=p^2$.
Let $q=p^a$ and $L=F^{1/q}$, and let $\varphi:L\xrightarrow{\sim}F$ be $x\mapsto x^q$.
Then
\[
 N_{L/F}=\varphi_*:K_2(L)\xrightarrow{\sim}K_2(F).
\]
Moreover, for every central division algebra $D$ over $F$, the reduced norm gives an isomorphism
\[
 \Nrd_2:I_F(D)_{(p)}\xrightarrow{\sim}K_2(F)_{(p)}.
\]
\end{lemma}
\begin{proof}
We use Matsumoto's identification $K_2=K_2^M$.
Since $[L:F]=q^2$, one has $N_{L/F}\res_{L/F}=q^2$, and on symbols one also has $\varphi_*\res_{L/F}=q^2$.
Since symbols generate $K_2(L)$ and each $q^2\{u,v\}=\{u^q,v^q\}$ lies in the image of restriction, it follows that $q^2(N_{L/F}-\varphi_*)=0$.
Since $K_2(F)$ has no $p$-torsion, we have $N_{L/F}=\varphi_*$ \cite{Matsumoto}, \cite[\S2.4, Theorem 6]{Izhboldin2000}.

First suppose that $D$ is $p$-primary, and choose $q$ divisible by its period.
Pullback along the $q$th-power Frobenius acts on $\Br(F)$ as multiplication by $q$, since it raises the defining $2$-cocycles to the $q$th power.
Since the composite $F\subset L\xrightarrow{\varphi}F$ is $x\mapsto x^q$, we have $\varphi_*[D_L]=q[D]=0$ in $\Br(F)$, and $D_L$ splits because $\varphi$ is an isomorphism.
Since the definition of $I_F(D)$ includes finite purely inseparable splitting fields, the splitting transfer $\theta_L$ is defined.

For any finite extension $F'/F$, the tower formula and $[{F'}^p:F^p]=[F':F]$ give $[F':{F'}^p]=p^2$, so the first assertion, applied to $F'$, shows that $N_{{F'}^{1/q}/F'}$ is an isomorphism.
Since $F'L\subset{F'}^{1/q}$, this norm is the composite $N_{F'L/F'}N_{{F'}^{1/q}/F'L}$, so $N_{F'L/F'}$ is surjective.
In particular, for every finite splitting field $E/F$ and $z\in K_2(E)$, write $z=N_{EL/E}w$.
Transitivity of splitting transfers gives
\[
 \theta_Ez=\theta_{EL}w=\theta_LN_{EL/L}w.
\]
Consequently $I_F(D)=\theta_LK_2(L)$, and $\Nrd_2$ maps it isomorphically onto $K_2(F)$ since $\Nrd_2\theta_L=N_{L/F}$ is an isomorphism.

For general $D$, split the prime-to-$p$ Brauer components over a finite extension $F'/F$ of degree prime to $p$.
Restriction and transfer preserve the corresponding splitting transfer subgroups, and the composite of these two maps is multiplication by this degree.
We use $\Nrd_2\circ\cor_{F'/F}=N_{F'/F}\circ\Nrd_2$ \cite[Proposition 6.6.1(1)]{KahnLevine2010}.
For surjectivity, given $z\in K_2(F)_{(p)}$, the transfer of a preimage of $\res_{F'/F}z$ over $F'$ is a preimage of $[F':F]z$, and $[F':F]$ is a unit in $\Z_{(p)}$.
For injectivity, restriction is injective on $I_F(D)_{(p)}$ for the same reason, so injectivity over $F'$ suffices.
\end{proof}
